\chapter{The Vendor Map}\label{ch:nw:the-vendor-map}

Eurex's list of \omterm{def:nw:the-vendor-map:isv}{independent software vendors} runs to thirty-three companies, each marked by the part of a trading firm it
serves: back office, middle office, front office. Its other lists name market data vendors, brokers and sponsored access
providers. A new trading firm buys most of its stack, and its first map of the market is a map of vendors: whom it buys
\omterm{def:nw:colocation-products-and-how-they-are-sold:colo}{colocation} from, which \omterm{def:nw:the-vendor-map:tickerplant}{ticker plant} decodes its feeds, whose servers run its strategies, whose clocks stamp its records.

The earlier chapters met these vendors one product at a time: switches and layer-1 devices in chapter~2, network cards in
chapter~3, clocks in chapter~4, capture in chapter~5, FPGA cards in chapters~6 and~7, servers in chapter~8, \omterm{def:nw:colocation-products-and-how-they-are-sold:colo}{colocation} and
routes in chapters~9 to~15. This chapter puts them on one map by category, attaches a firm's stack to it as a graph, and
measures what the map shows: how concentrated the firm's spend is and which vendors' failure would stop it trading. Vendors
are named only from a venue's list or their own publications; the example stack is synthetic.

\section{Market data}

\begin{definition}[Ticker plant]\label{def:nw:the-vendor-map:tickerplant}
A \emph{ticker plant}\index{ticker plant} is a system, bought or built, that receives a firm's market data feeds once, decodes
and normalises them into one format, maintains the order books, and distributes the results to the firm's trading, risk and
monitoring applications over its internal network.
\end{definition}

A firm can decode each feed inside each strategy (Book~13's feed handler) or once, centrally, in a \omterm{def:nw:the-vendor-map:tickerplant}{ticker plant}.
Vendors sell both. Exegy describes centralised \omterm{def:nw:the-vendor-map:tickerplant}{ticker plants} as ``the dominant architectural response to the scaling
limitations of embedded feed handlers'', and names their cost: ``distribution complexity and latency overhead''
(\cref{dat:nw:the-vendor-map:lists}). Pico's product page defines the \omterm{def:nw:the-vendor-map:tickerplant}{ticker plant} as software that connects to the feeds,
normalises the messages and keeps the books in memory, for more than 230 venues. The choice is chapter~1's multicast question
turned into procurement: one decode and a second network, or many decodes and none.

\section{Execution software and order management}

\begin{definition}[Independent software vendor]\label{def:nw:the-vendor-map:isv}
An \emph{independent software vendor}\index{independent software vendor} is a firm that sells trading, order management,
risk or post-trade software to trading firms and that an exchange certifies to connect to its systems on their behalf, without
itself being a member.
\end{definition}

Exchanges publish their certified vendors. Eurex's list gives each company's services by office: Trading Technologies for back
and front office, Broadridge and OnixS for front office among many others. For a new firm the list is a shopping list and a
constraint: a vendor that is not certified for a venue cannot connect the firm to it, and every venue's certification is a
separate project (Book~13).

\begin{dated}{2026-09}{Vendors on the map, from venues' lists and vendors' own publications}\label{dat:nw:the-vendor-map:lists}
\begin{itemize}
\item Eurex \omterm{def:nw:the-vendor-map:isv}{independent software vendors}: 33 companies, each with its offices served, among them Trading Technologies (back
  and front office), Broadridge (front office), OnixS (front office).
\item Market data: Exegy (centralised \omterm{def:nw:the-vendor-map:tickerplant}{ticker plants}, FPGA feed handling); Pico (InRush Ticker Plant, more than 230 venues).
\item Timing: Meinberg (IEEE 1588 \omterm{def:nw:time-synchronisation:gm}{grandmaster clocks}, LANTIME time servers). Monitoring: Pico (Corvil Analytics).
\item From earlier chapters: switches (Cisco, Arista), network cards and FPGA cards (AMD), overclocked servers (Blackcore,
  Business Systems International), \omterm{def:nw:colocation-products-and-how-they-are-sold:colo}{colocation} (Equinix), a microwave route (Quincy Data).
\end{itemize}
\end{dated}

\section{Network hardware and servers}

The hardware categories are the ones measured in chapters~2 to~8, and each has a handful of vendors that low-latency firms use: a
switch vendor for the aggregation layer and a layer-1 device vendor for fan-out; a network card that supports kernel bypass; an
FPGA card for the hardware path; overclocked servers. A firm usually standardises on one vendor per category, for spares, skills
and drivers; that choice is what makes a vendor a single point of failure, and chapter~28's recovery plan has to cover it.

\section{Timing and monitoring}

Timing and monitoring are the categories a firm can survive losing for an hour and must not lose for a day. A \omterm{def:nw:time-synchronisation:gm}{grandmaster clock}
feeds every timestamp the regulators ask for (chapter~4); network analytics and \omterm{def:nw:capturing-and-monitoring-the-network:capture}{capture appliances} (chapter~5) see what the firm's
own systems do not. Neither sits on the path from data to orders, so neither stops trading; both stop the firm from proving what it
did, which is its own kind of halt.

\section{Reading the map: lock-in, concentration and exit}

A stack is a graph: market data comes in, passes through components, and leaves as orders. Each component is supplied by one or
more vendors, or built in house. Two questions follow. Which components, failing alone, cut every path from data to orders? And
which vendors, failing, take out every component they alone supply and so cut it?

\omcode{code/firm/vendormap/firm_vendormap.py}{60}{83}{Paths from data to orders, the components that cut them alone, and the vendors that do.}\label{lst:nw:the-vendor-map:graph}

\begin{proposition}[Vendor points of failure]\label{prop:nw:the-vendor-map:vendors}
A vendor halts trading if and only if the components it alone supplies form a cut between data and orders; a component supplied
by two vendors is never removed by one vendor's failure, so dual-sourcing a component removes it from every vendor's cut, but not
from its own.
\end{proposition}

\begin{proof}
A vendor's failure removes exactly the components that have no other supplier; trading halts if and only if no path from data to
orders survives their removal. A dual-sourced component survives either vendor's failure, but the component itself, failing for
another cause (its software, its configuration), still cuts the path if it is the only route.
\end{proof}

The chapter's example firm trades from Secaucus. Its stack: \omterm{def:nw:colocation-products-and-how-they-are-sold:colo}{colocation} and \omterm{def:nw:colocation-products-and-how-they-are-sold:mmr}{cross-connects} from Equinix; a \omterm{def:nw:switches-and-layer-1-devices:l1}{layer-1 fan-out} from Arista
and switches from Cisco, in parallel; an Exegy \omterm{def:nw:the-vendor-map:tickerplant}{ticker plant} and in-house feed handlers, in parallel; Blackcore servers with AMD network
cards; its own execution gateway; and, off the path, order management from Trading Technologies, Meinberg grandmasters, Pico network
analytics and a Quincy Data microwave route. It spends USD 3.19 million a year on them (synthetic figures).

\begin{omfigure}
\begin{tikzpicture}[font=\footnotesize,node distance=4mm,x=1.13cm,
  spof/.style={draw=omThm,fill=omThm!25,rounded corners,align=center,minimum height=8mm,inner sep=2pt,text width=2.3cm},
  ok/.style={draw=omDef,fill=omDef!10,rounded corners,align=center,minimum height=8mm,inner sep=2pt,text width=2.3cm},
  off/.style={draw=gray,fill=gray!8,rounded corners,align=center,minimum height=8mm,inner sep=2pt,text width=2.3cm,dashed}]
\node[spof] (colo) at (0,0) {colocation\\Equinix};
\node[spof] (xc) at (2.6,0) {cross-connects\\Equinix};
\node[ok] (l1) at (5.2,0.75) {layer-1 fan-out\\Arista};
\node[ok] (sw) at (5.2,-0.75) {switches\\Cisco};
\node[ok] (tp) at (7.8,0.75) {ticker plant\\Exegy};
\node[ok] (fh) at (7.8,-0.75) {feed handlers\\in house};
\node[spof] (srv) at (10.4,0) {servers\\Blackcore};
\node[spof] (nic) at (10.4,-2.0) {network cards\\AMD};
\node[spof] (gw) at (7.8,-2.6) {gateway\\in house};
\node[off] (oms) at (0,-2.2) {orders (OMS)\\Trading Technologies};
\node[off] (clk) at (2.6,-2.2) {grandmasters\\Meinberg};
\node[off] (mon) at (5.2,-2.6) {analytics\\Pico};
\draw[->] (colo) -- (xc); \draw[->] (xc) -- (l1); \draw[->] (xc) -- (sw);
\draw[->] (l1) -- (tp); \draw[->] (sw) -- (fh); \draw[->] (l1) -- (fh); \draw[->] (sw) -- (tp);
\draw[->] (tp) -- (srv); \draw[->] (fh) -- (srv); \draw[->] (srv) -- (nic); \draw[->] (nic) -- (gw);
\end{tikzpicture}

\omcaption{The example firm's stack from data (left) to orders (the gateway). Filled red: components whose failure alone stops
trading; blue: components with a parallel alternative; dashed grey: off the trading path. Synthetic stack; vendors as in
\cref{dat:nw:the-vendor-map:lists}. Data: \texttt{nw\_vendors.report()}.}\label{fig:nw:the-vendor-map:stack}
\end{omfigure}

Five components cut the path alone: \omterm{def:nw:colocation-products-and-how-they-are-sold:colo}{colocation}, \omterm{def:nw:colocation-products-and-how-they-are-sold:mmr}{cross-connects}, the strategy servers, their network cards and the in-house gateway.
Three vendors do: AMD, Blackcore and Equinix; Arista, Cisco and Exegy do not, because each has a parallel path. The Herfindahl--Hirschman
index of the firm's spend by vendor (shares of total spend, in-house work excluded from the index) is 1\,149 on a scale of zero to
10\,000, which reads as spread; but the index measures spend, not dependence: the
firm's largest vendor by spend, Equinix at 22.6\%, is a point of failure, and its second, Quincy Data at 15.0\%, is not.

\begin{omfigure}
\begin{tikzpicture}
\begin{axis}[xbar,width=10.5cm,height=6.4cm,bar width=6pt,xmin=0,xmax=27,ytick=data,y dir=reverse,
  yticklabels from table={figdata/networks/27-the-vendor-map/shares.csv}{vendor},
  xlabel={share of yearly spend (\%)},nodes near coords,every node near coord/.append style={font=\scriptsize},
  nodes near coords={\pgfmathprintnumber[fixed,fixed zerofill,precision=1]{\pgfplotspointmeta}},
  tick label style={font=\footnotesize},label style={font=\small},grid=major,grid style={gray!25},table/col sep=comma]
\addplot[fill=omDef!55,draw=omDef] table[y=x,x=share]{figdata/networks/27-the-vendor-map/shares.csv};
\end{axis}
\end{tikzpicture}

\omcaption{The example firm's yearly spend by vendor (synthetic figures): the spend is spread, the dependence is not. Data:
\texttt{fig\_vendors.py}, \texttt{nw\_vendors.report()}.}\label{fig:nw:the-vendor-map:shares}
\end{omfigure}

Adding a second server vendor on half the servers lowers the index to 1\,119 and removes Blackcore from the vendor points; the strategy
servers remain a single component that can fail for other reasons. The exit times of the five single-point components add to 48 weeks:
leaving Equinix alone is half a year. That is the practical meaning of vendor lock-in (Book~16, chapter~20, treats it as a business
decision): not a contract clause but the weeks it takes to replace a component on the trading path.

\begin{method}[Reading a firm's vendor map]\label{met:nw:the-vendor-map:read}
\begin{enumerate}
\item List every component on the path from market data to orders and every one off it; attach each to its vendors, with the source
  that names the vendor and the date.
\item Draw the graph; find the single components and the vendors whose failure cuts it.
\item Compute the concentration of spend, and read it beside the vendor points: they answer different questions.
\item For each vendor point, estimate the exit time and decide: dual-source, keep spares, or accept and cover in the recovery plan
  (chapter~28).
\item Export the bill of materials to the connectivity plan (chapter~29) and re-run the map at every contract renewal.
\end{enumerate}
\end{method}

\section{Tutorial: one vendor too many}\label{tut:nw:the-vendor-map:tut}

\begin{tutorial}
\textbf{Goal.} Build the vendor map as data, attach a stack, and find its single points of vendor failure.
\textbf{End state:} \cref{fig:nw:the-vendor-map:stack,fig:nw:the-vendor-map:shares}.
\begin{tutsteps}
\item \textbf{Inventory.} \texttt{firm\_vendormap.load\_inventory} reads the dated vendor rows, each with its ledger source.
\item \textbf{Stack.} \texttt{nw\_vendors.COMPONENTS} and \texttt{EDGES} describe the example firm; \texttt{Stack} holds them.
\item \textbf{Measures.} \texttt{single\_points}, \texttt{vendor\_points} (\cref{lst:nw:the-vendor-map:graph}), \texttt{hhi} and
  \texttt{bom}; \texttt{second\_server\_vendor} for the change.
\end{tutsteps}
\textbf{What to change next.} Weight each vendor point by its exit time and the probability of its failure, and rank the fixes by
risk removed per dollar.
\end{tutorial}

\section{Build: the vendor map}\label{bld:nw:the-vendor-map:vendormap}

\begin{build}
\bfield{Purpose} A vendor inventory and a firm's stack as a graph, with concentration, single points of failure and a bill of
materials: the vendor rows of chapter~29's plan.
\bfield{Interface} \texttt{firm\_vendormap}: \texttt{InventoryRow}, \texttt{load\_inventory}, \texttt{Component}, \texttt{Stack}
(\texttt{connected}, \texttt{single\_points}, \texttt{vendor\_points}, \texttt{shares}, \texttt{hhi}, \texttt{bom}).
\bfield{Rules} Every vendor row names its ledger source and date; the stack and its spend are the firm's own (here synthetic).
\bfield{Acceptance tests} \texttt{code/firm/vendormap/tests/}: inventory sources, a three-component graph by hand, and in-house
components excluded from vendor points.
\bfield{Stretch} Failure probabilities and exit costs; shared sub-suppliers (two vendors on one chip).
\end{build}

\begin{omsources}
\item Eurex, participant lists: independent software vendors.
\item Exegy, \emph{Design Patterns for Market Data, Part~2: Centralized Ticker Plant}; Pico, InRush Ticker Plant; Meinberg,
  product range.
\item This book's chapters~2 to~14 for the hardware, colocation and route vendors.
\end{omsources}

\section{Exercises}

\begin{exercise}[$\star$]\label{exo:nw:the-vendor-map:1}
What does a \omterm{def:nw:the-vendor-map:tickerplant}{ticker plant} do, and what does centralising it cost?
\end{exercise}

\begin{exercise}[$\star$]\label{exo:nw:the-vendor-map:2}
What is an \omterm{def:nw:the-vendor-map:isv}{independent software vendor}, and why does an exchange certify it?
\end{exercise}

\begin{exercise}[$\star$]\label{exo:nw:the-vendor-map:3}
Why do the \omterm{def:nw:time-synchronisation:gm}{grandmaster clocks} not appear among the example firm's points of failure, and why should the firm still care?
\end{exercise}

\begin{exercise}[$\star\star$]\label{exo:nw:the-vendor-map:4}
Using \cref{prop:nw:the-vendor-map:vendors}, why is Exegy not a vendor point of failure for the example firm?
\end{exercise}

\begin{exercise}[$\star\star$]\label{exo:nw:the-vendor-map:5}
Compute the example firm's index with Equinix's spend doubled.
\end{exercise}

\begin{exercise}[$\star\star$]\label{exo:nw:the-vendor-map:6}
The firm's second server vendor buys its network cards from the same maker. What does the map miss?
\end{exercise}

\begin{exercise}[$\star\star\star$]\label{exo:nw:the-vendor-map:7}
\emph{Coding.} With \texttt{firm\_vendormap.Stack}, add a second \omterm{def:nw:colocation-products-and-how-they-are-sold:colo}{colocation} site, from another operator, with its own \omterm{def:nw:colocation-products-and-how-they-are-sold:mmr}{cross-connects} feeding the same servers.
Which vendor points remain?
\end{exercise}

\begin{exercise}[$\star\star\star$]\label{exo:nw:the-vendor-map:8}
\emph{Find the flaw.} ``Our largest vendor is only 23\% of our spend, so we are not dependent on anyone.''
\end{exercise}

\section{Problem: One Vendor Too Many}

\begin{problem}[{Weekend problem --- concentration and single points of vendor failure}]\label{pb:nw:the-vendor-map:1}
A new firm trades from Secaucus with the chapter's example stack. Its board asks which vendors could stop it trading and what it
would take to change that.

\medskip\noindent\textbf{Part I --- The map.}
\begin{enumerate}
\item Which categories does the firm buy, and from whom?
\item What sources name each vendor?
\item What does the firm spend in a year, and on whom most?
\item What is the concentration index?
\end{enumerate}

\medskip\noindent\textbf{Part II --- The graph.}
\begin{enumerate}[resume]
\item Which components cut the path alone?
\item Which vendors do?
\item Why are Arista, Cisco and Exegy not among them?
\item What are the exit times of the single-point components?
\end{enumerate}

\medskip\noindent\textbf{Part III --- The changes.}
\begin{enumerate}[resume]
\item What does a second server vendor change?
\item What would remove AMD from the vendor points?
\item Why can Equinix not be removed cheaply?
\item What does the off-path software do to the firm if it fails?
\end{enumerate}

\medskip\noindent\textbf{Part IV --- The verdict.}
\begin{enumerate}[resume]
\item State the \emph{named result}: the concentration index of the firm's stack and the list of components whose single vendor's
  failure halts trading.
\item Which inputs are published and which the firm's own?
\item Why do the index and the vendor points disagree?
\item What should go into the recovery plan of chapter~28?
\item What should the board ask at each contract renewal?
\item How does the map change for a firm that builds its own \omterm{def:nw:the-vendor-map:tickerplant}{ticker plant}?
\item Where does this go in chapter~29's plan?
\item In one sentence: what makes a vendor dangerous?
\end{enumerate}
\end{problem}

\section{Interview questions}

\begin{interviewq}[$\star$ \iqroles{developer}]\label{iq:nw:the-vendor-map:1}
\omterm{def:nw:the-vendor-map:tickerplant}{Ticker plant} or embedded feed handlers: what are the trade-offs?
\end{interviewq}

\begin{interviewq}[$\star\star$ \iqroles{developer}]\label{iq:nw:the-vendor-map:2}
How would you find the single points of failure in a trading firm's infrastructure?
\end{interviewq}

\begin{interviewq}[$\star\star$ \iqroles{developer, researcher}]\label{iq:nw:the-vendor-map:3}
Your \omterm{def:nw:the-vendor-map:tickerplant}{ticker plant} vendor announces the end of support for your version in six months. What do you do?
\end{interviewq}

\begin{interviewq}[$\star\star$ \iqroles{developer}]\label{iq:nw:the-vendor-map:4}
Why would a firm standardise on one server vendor, and when should it not?
\end{interviewq}

\begin{interviewq}[$\star\star$ \iqroles{trader}]\label{iq:nw:the-vendor-map:5}
Your order management vendor is down for an hour. What still works, and what do you do?
\end{interviewq}

\begin{interviewq}[$\star\star\star$ \iqroles{developer, researcher}]\label{iq:nw:the-vendor-map:6}
Design the vendor strategy of a new trading firm for its first two years: what to buy, what to build, and how to keep the exits open.
\end{interviewq}
